\documentclass{article}
\usepackage{graphicx} % Required for inserting images
\usepackage{amsmath}
\usepackage{amssymb}
\title{Topology}
\author{Maxwel Kim}
\date{June 2026}

\begin{document}

\maketitle

\section*{Euclidian set}

\textbf{Archemedis Theorem} states that 

\[
1. \forall x>0, y\in \mathbb{R}, \exists n \; \text{such that} \; nx>y
\]

\[
2. \forall x>0, \exists n \; \text{such that} 0<\frac{1}{n}<x
\]

\[
3. \forall n\in \mathbb{N} \; \exists x=0 \; \text{such that} \; 0\leq x<\frac{1}{n} 
\]

by proof by contradiction, let's say

\[
nx\leq y
\]

so that by supreme theorem

\[
\exists \; \alpha =\sup A(\in \mathbb{R})
\]

but since $x>0$, there exists $m\in \mathbb{N}$ such that $\alpha-x<\alpha$ satisfies. therefore since $\alpha-x$ is not a supreme, it comes to $\alpha-x<mx$ so that $m$ exists, but rearranging, $\alpha<(m+1)x\in A$ makes it that $\alpha$ is not a supreme of $A$ therefore it works out.

\textbf{Open set} is said when 

\[
\forall x\in U, \exists x\in B(x,r)\subset U
\]

so that $r$ exists. and $B(x,r)$ is an open ball of $\mathbb{R}^n$ 

On the other hand, it is called \textbf{closed set} if 

\[
C^c=\mathbb{R}^n-C
\]

applies so that $C^c$ is an open set that makes $C$ a closed set

here is an example, statement "every finite subset is closed set in $\mathbb{R}^n$ is true because

\[
\forall A=\{a_1,....,a_k\}\subset \mathbb{R}^n
\]

then we have to prove that 

\[
A^c=\mathbb{R}^n-A
\]

is an open set, so we pick a point $p\in A^c$ and now can state

\[
r := \min\{d(p,a_i)|i\in \{1,2,...,k\}\}
\]

then since $B(p,r)\cup A=\varnothing$ we can then claim that $p\in B(p,r)\subset A^c$ so $A^c$ is open set of $\mathbb{R}^n$ A is closed set of $\mathbb{R}^n$

now comes definition of continuity.

\textbf{Definition:} continuous at point $p$ occurs when for function

\[
f:\mathbb{R}^m \rightarrow \mathbb{R}^n
\]

we can claim 

\[
\forall \epsilon>0, \;d(x,p)<\delta \rightarrow d(f(x),f(p))<\epsilon, \exists \delta>0
\]

\textbf{Topology} is when

\[
\varnothing, X\in \mathcal{T} \; \forall X \text{ being set}
\]
\[
\forall U_\alpha\in\mathcal{T}(\alpha\in \Lambda), \; \bigcup_{\alpha\in \Lambda}U_\alpha \in \mathcal{T}
\]
\[
\forall \text{finite } U_\alpha \in \mathcal{T},(\alpha\in\{1,2,...,m\}), \bigcap_{\alpha=1}^{m} U_\alpha\in \mathcal{T}
\]

all applies

and of course, when there is toplogy space $(X,\mathcal{T})$, and $U\in \mathcal{T}$, it applies that $U$ is \textbf{$\mathcal{T}-$open set}

\textbf{for example,} say we have a set of 

\[
\mathcal{T}_1=\{\varnothing, X, \{a,b\}, \{b,c\}\}
\]

this is not a topological space because intersection of $\{a,b\}, \{b,c\}$ is $\{b\}$ and doesn't exist in the set, so it would be topological space if we add $\{b\}$ in there.

\textbf{Subspace} is when

\[
\mathcal{T}_A=\{A\cap U| U\in \mathcal{T}\} \forall A\subset X
\]

and we call $(A,\mathcal{T}_A)$ is subspace of $(X,\mathcal{T})$

\textbf{Neighborhood} is defined as 

\[
\forall x\in U, U\in \mathcal{T}, \; x\in U \subset N 
\]

validates, then it's called $N$ is neighborhood of $x$.

now definition for closure,

\textbf{closure} is when defined as

\[
\bar{A} = \cap \{C|C \; \text{is a closure of } A\subset C\}
\]

and here follows the theorem:

\[
A\subset \bar{A} \subset C
\]

if $C$ includes $A$ as closed set. Also to be able to be closed set, it has to be

\[
A=\bar{A}
\]

for an obvious reason from definition of closure.

another self-explanatory expression is this:

\[
\bar{A}=A\cup A'
\]

if A is an closed set, then boundary of $A$ will be included in $A$ which makes sense as closure, if it's open set, then $A'$(boundary) will be the missing piece from open to closed set.

\section*{Basis}

Basis occurs when

\[
\forall (X,\mathcal{T}), \mathcal{B}\subset \mathcal{T} 
\]

and to be able to be a basis, it has to have these requirements:

\[
X=\bigcup_{\lambda\in\Lambda}B_\lambda, where \; \; B_\lambda\in\mathcal{B}
\]
\[
\forall B_1, B_2\in \mathcal{B}, B_1\cap B_2 \in \mathcal{B}
\]

\textbf{Sub-basis} is defined as

\[
\forall S=\{S_\alpha|S_\alpha\subset X, \alpha\in \Lambda\}, \bigcup S_\alpha=X
\]

happens to be true

for $X=(X,\mathcal{T}_1), Y=(Y,\mathcal{T}_2)$ \textbf{Continuity} is defined as function $f:X\rightarrow Y$ being

\[
\forall f(x)\in V \text{ where } x \text{ is in X}, V\in \mathcal{T}_2
\]

\[
f(U)\subset V, \exists U
\]

therefore as defined as 

\[
f:X\rightarrow Y, \forall V\in \mathcal{T}_2 f^{-1}(V)\in \mathcal{T}_1
\]

given $X=\{x,y,z\}, \mathcal{T}_1=\{X,\varnothing,\{x,y\},\{z\}\}, Y=\{a,b\}, \text{and } \mathcal{T}_2=\{Y,\varnothing, \{b\}\}$

we define function $f$ as this:

\[
f:X\rightarrow Y \; \textbf{ s.t.} \; f(x)=b, f(y)=a, f(z)=b
\]

and this is clearly not a continuous map because 

\[
f^{-1}(\{b\})=\{x,z\}\notin \mathcal{T}_1
\]

this leads to the fact that

\textbf{Theorem} $f$ being continuous throughout all point at domain $X$ will mean same as being continuous map

\textbf{Definition:} Homeomorphism is topology space's function that appears to apply these:

\[
f \; \text{ bijective}
\]
\[
f \text{ continuous}
\]
\[
f^{-1} \text{continuous}
\]

when these 3 applies it appears to be homeomorphic.

for example, function from topology $(X,\mathcal{U}_X) , (Y,\mathcal{U}_Y)$ is given and for

\[
X=\mathbb{Z}, Y=\{\frac{1}{n}|n\in \mathbb{Z}-\{0\}\}\cup\{0\}
\]

this shows that it is not homeomorphic. this is because if we say $\{0\}\in \mathcal{U}_Y$ then there has to be $\mathcal{U}$ that includes $U\cap Y=\{0\}$. therefore, there exists $0\in (a,b)$ such that both $(a,b)\subset U$ and $Y\cap (a,b)\subset Y\cap U=\{0\}$. but $0\in (a,b)$ doesnt exist because $Y\cap (a,b)$ is an infinite set. therefore this is not homeomorphism

\textbf{theorem}

\[
f \text{ is homeomorphism}
\]
\[
f \text{ is continuous and open map}
\]
\[
f \text{ is continuous and closed map}
\]
\[
\forall A\subset X, \; \bar{f(A)}=f(\bar{A})
\]

these 4 are all equivalent statement.


\textbf{Distance function}

\textbf{metric function } is defined as function $d:X \times X \rightarrow \mathbb{R}$ when

\[
d(x,y)\geq 0
\]
\[
d(x,y)=0 \longleftrightarrow x=y
\]
\[
d(x,y)=d(y,x)
\]
\[
d(x,z)\leq d(x,y)+d(y,z)
\]

and here is definition for open ball:

\[
B_d(p,r)=\{x\in X| d(p,x)<r\}
\]

and for to be able to define continuity,

\[
f(B_d(x,\delta))\subset B_\rho(f(x),\epsilon)
\]

\section*{Spaces}

it's called sequence if 

\[
f:\mathbb{N} \rightarrow X, f(n)=x_n \in X
\]

\[
(x_n)_{n\in \mathbb{N}}:=<x_n>
\]


\textbf{ theorem:} if $f:X\rightarrow Y$ is continuous, it is sequentially continuous

this is because if we let $<x_n>$ as one sequence in $X$, and $\forall p\in X$ we say $<x_n>\rightarrow p$ then  we need to prove $<f(x_n)> \rightarrow f(p)$ to prove that it's continuous. since $f$ is continuous, $f(p)\in Y$ occurs and $f(U)\subset V$ also occurs so that $p$ exists in $U$ therefore $f(x_n)\in f(U)\subset V$ this leads to $<f(x_n)> \rightarrow f(p)$, therefore it is sequentially continuous.

it is called \textbf{ first countable space} if for all $x$ where its countable local base at $x$ exists in $X$.

example: 

if $X$ is incountable set, then cofinite space $(X,\mathcal{T}_f)$ is NOT first countable space.

this is because if we set $(x,\mathcal{T}_f)$ as first countable space, then this leads to the fact that

\[
\forall p\in X, \mathcal{B}_p=\{B_n|n\in \mathbb{N}\} 
\]

exists. Then for every $n\in \mathbb{N}$, we can say $p\in B_n\neq \varnothing$ and since $B_n\in \mathcal{T}_f$, $B_n^c$ is therefore finite set, then

\[
(\bigcap_{n\in \mathbb{N}}B_n)^{c}=\bigcup_{n\in\mathbb{N}}B_n^c
\]

is countable set. and since $X$ is not countable, the big cup is also not countable. so there exists point where

\[
p\neq q, q\in \bigcap_{n\in \mathbb{N}}B_n
\]

so this is not valid.

\textbf{theorem:}

first countable space has hereditary property\\
first countable space has countable productive property\\

we need to prove that its subspace also is first countable space to be hereditary.

\[
(A,\mathcal{T}_A)\subset (X,\mathcal{T})
\]
\[
\forall a\in A\subset X, \{B_n\cap A|n\in \mathbb{N}\}\in \mathcal{T}_A
\]

and furthermore

\[
a\in B_n\cap A
\]

so it is proven to be hereditary

for countable finite product property:

\[
\forall x=(x_1,...,x_n)\in \prod_{\alpha\in \Lambda}X_\alpha 
\]

then 

\[
\mathcal{B}=\{B_1\times B_2 \times ... \times B_n|B_\alpha \in \mathcal{B}_{x_\alpha}, \alpha\in \Lambda\}
\]

and this point $x\in\prod X_\alpha$ is countable, therefore is first countable space

for infinite countable product property it's different:

\[
\forall x=(x_\alpha)_{\alpha \in \mathbb{N}}\in \prod_{\alpha\in \mathbb{N}}X_\alpha 
\]

then since it's countable for every $\alpha$, $\exists \mathcal{B}_{x_\alpha}$ countable local basis where $B_\alpha=X_\alpha$ except finite $\alpha$ therefore it is first countable space

\textbf{theorem:} if first countable space $X$, $f:X\rightarrow Y$ is sequentially continuous, then it is continuous.

\textbf{definition: } second countable space is defined when for $(X,\mathcal{T})$ exists countable basis. 
\\


second countable space is first countable space because by definition, first countable space is countable local basis whereas second countable space consists of countable base.

for similar reason, second countable space has hereditary property and countable product property. 

here's the different between those two space:
\\


second countable space has topological property:

\[
f:(X,\mathcal{T}_1)\rightarrow (Y,\mathcal{T}_2)
\]

say this function is homeomorphism, and domain is second countable space, then

\[
\exists \mathcal{B}=\{B_\alpha|\alpha\in \mathbb{N}\}
\]

where

\[
\{f(B_\alpha)|\alpha\in\mathbb{N}\}
\]

is range's countable basis.

\section*{Separation Axiom}

$T_0$ space also known as Kolmogorov space is defined as 

\[
\forall a,b\in X, \exists \; U\in \mathcal{T} \text{ s.t. } U(\ni a), U(\not\ni b) \; \text{ or } U(\not\ni a) , U(\ni b)
\]

this is a theorem of definition-based

\[
1. X \text{ is } T_0 \text{ space}
\]
\[
2. \forall a,b\in X \text{ applies that } a \notin \bar{\{b\}} \text{ or } b \notin \bar{\{a\}}
\]

now what this means is that 

\[
\forall a,b\in X, 
\]

as definition of $T_0$ space, 

\[
\exists U \; \text{such that } a\in U, b\notin U \text{ or } a\notin U, b\in U. 
\]

so for the first case,

\[
b\in U^c(\text{closed set)}, \bar{\{b\}}\subset U^c. \rightarrow a\notin \bar{\{b\}}
\]

second case works the same.

and vice versa from $2\rightarrow 1$ works the same other way around.\\

easy to prove this space has product property, hereditary property and topological property.\\

$T_1 $ Space is defined as 

\[
\forall U,V(\text{opened set}) \text{ such that } a\in U, b\in V \text{ & } a\notin V, b\notin U
\]

also has product property, hereditary property and topological property.\\

\textbf{theorem:} there two statements are equivalent

\[
1. \text{ topology } X \text{ is } T_1 \text{ space}
\]
\[
2. \forall a\in X, \{a\}= \text{closed set}
\]

this is because 

\[
\forall a\in V^c, \{a\}\subset V^c \text{ if we keep the definition.} 
\]
\[
b\in V\subset \{a\}^c
\]

therefore open set is $\{a\}^c$ and $\{b\}^c$ and $U=\{b\}^c$ and $V=\{a\}^c$ therefore $X $ is $T_1$ space.

$T_1$ space also has product property, hereditary property and topological property

for proving this, let's say we have subspace $(A,\mathcal{T}_A)$ that's $T_1$ space. then for $a,b\in A$, they are in $a,b \in X$ since they are subspace of $X$, and since $X$ is $T_1$ space all the definition stuff works on $a,b$. 

\[
a\in U_A, b\in V_A, a\notin V_A, b\notin U_A
\]

therefore it is $T_1$.

for topological property, let's say $(X,\mathcal{T}_1)$ and $(Y,\mathcal{T}_2)$ are homeomorphism and $(x,\mathcal{T}_1)$ $T_1$ space.

then there exists homeomorphic function $h:(X,\mathcal{T}_1)\rightarrow (Y,\mathcal{T}_2)$ now we have to prove topology $Y$ to be $T_1$ space.

To do this, we have $y_1,y_2\in Y$ such that since it's homeomorphism 

\[
h^{-1}(y_1)=x_1, h^{-1}(y_2)=x_2
\]

and since $X$ space is homeomorphic, we can say

\[
x_1\in U_1\in \mathcal{T}_1, \; x_2\in U_2 \in \mathcal{T}_2
\]

setting up $Y$ space as $V$, it happens the same where

\[
y_1\in V_1, y_2\in V_2
\]

therefore topological space works.


\textbf{Hausdorff space} or $T_2$ space is defined as 

\[
\forall a\in U, b\in V \; \exists U\cap V=\varnothing
\]

to prove $T_2$ space being hereditary, 

\[
\forall (A,\mathcal{T}_A)\subset (X,\mathcal{T})
\]

then

\[
\forall a,b\in A, a,b\in X \rightarrow a\in U\in \mathcal{T}, b\in V\in \mathcal{T}
\]

such that

\[
U\cap V=\varnothing
\]

therefore

\[
a\in A\cap U:=U_A\in \mathcal{T}_A, b\in A\cap V:= V_A\in \mathcal{T}_A
\]

therefore it is hausdorff space.

to prove it's topologically proven, we set 2 different homeomorphic topology and set the domain as $T_2$, like previous proof.

\[
(X,\mathcal{T}_1)(T_2)\rightarrow (Y,\mathcal{T}_2).
\]

then as always, we have to prove range is $T_2$. and to do that

\[
\forall y_1,y_2\in Y \text{ such that } h^{-1}(y_1,y_2)=(x_1,x_2)
\]

this shows that there is open space $U$, that

\[
x_1\in U_1\in \mathcal{T}_1, x_2\in U_2\in \mathcal{T}_1, U_1 \cap U_2=\varnothing
\]

and the fact that it is homeomorphic, where bijective property now brings in

\[
h(U_1)=V_1, h(U_2)=V_2
\]

then

\[
\forall V_1, V_2\in \mathcal{T}_2, y_1\in V_1, y_2\in V_2
\]

therefore, it is $T_2$ space. 

to prove its product property, 

\[
\forall x,y \;, \; x_{\alpha_0} \in U_{\alpha_0}, y_{\alpha_0} \in V_{\alpha_0}. U_{\alpha_0}\cap V_{\alpha_0}=\varnothing
\]

there exists openset $U,V$ that applies this and product space also exist where:

\[
P^{-1}_{\alpha_0}(U_{\alpha_0})=U, P^{-1}_{\alpha_0}(V_{\alpha_0})=V
\]

where

\[
P_{\alpha_0}: \prod_{\alpha\in \Lambda} X_\alpha \rightarrow X_{\alpha_0}
\]

exists. Therefore it is product property.

\textbf{regular space} is defined as

\[
F\subset U, x\in V \text{ & } U\cap V = \varnothing
\]

this is also called $T_3$ space.

it is easy to prove that $T_3$ space is $T_2$ space because

\[
\{p\}\subset U, q\in V, U\cap V=\varnothing
\]

where $\{p\}$ is treated as space.and it works to prove it's $T_2$ space therefore.

for same proof, it can be proven that they have heredarity, topological, and product property.

\textbf{Normal space} is defined when

\[
F_1\subset U, F_2\subset V, U\cap V=\varnothing
\]

\textbf{theorem:} it can be easily proven in the same way that $T_4$ is $T_3$. 

\[
F\subset U, \{x\}\subset V, U\cap V=\varnothing
\]


\section*{Compact Space}

\textbf{covering } is defined as 

\[
A\subset \bigcup_{U\in \mathcal{C}}U
\]

$\mathcal{C}_0$ is called subcovering when $\mathcal{C}_0\subset \mathcal{C}$ where $\mathcal{C}$ is covering

now that we know covering, \textbf{compact set} is when $(X,\mathcal{T})$ being topological space and $A\subset X$. where 

\[
\mathcal{C}=\{U_\lambda|\lambda\in \Lambda, U_\lambda\in \mathcal{T}\} 
\]

includes finite subcovering 

\[
\mathcal{C}_0=\{U_{\lambda_1},...,U_{\lambda_n}\}.
\]\\


\textbf{Heine Borel theorem}
\begin{theorem}
In the usual topological space $(\mathbb{R}, u)$, the bounded closed set $[a,b]$ is compact.
\end{theorem}

\begin{proof}
Let $\mathcal{C} = \{U_\alpha \mid \alpha \in \Lambda\}$ be an arbitrary open cover of $[a,b]$. Define
\[
A = \{x \in [a,b] \mid [a,x] \text{ is covered by a finite subcover of } \mathcal{C}\}.
\]
Then $A \subset [a,b]$ and $a \in A$, so $A \neq \emptyset$; also $A$ is bounded above by $b$. Hence
\[
\sup A = c
\]
exists for some $c \in [a,b]$.

Suppose $c < b$. Since $c \in U_\alpha$ for some $U_\alpha \in \mathcal{C}$, there exists $\varepsilon > 0$ such that
\[
(c-\varepsilon, c+\varepsilon) \subset U_\alpha, \qquad c+\varepsilon < b. 
\]
By definition of $c$, there exist $U_{\alpha_1}, \dots, U_{\alpha_n} \in \mathcal{C}$ such that
\[
\left[a, c - \tfrac{\varepsilon}{2}\right] \subset \bigcup_{i=1}^n U_{\alpha_i},
\]

\[
\left[a, c + \tfrac{\varepsilon}{2}\right] \subset \left(\bigcup_{i=1}^n U_{\alpha_i}\right) \cup U_\alpha.
\]
Thus $c + \tfrac{\varepsilon}{2} \in A$, contradicting $c = \sup A$. Therefore $c = b$, i.e.,
\[
[a,b] \subset U_{\alpha_1} \cup \cdots \cup U_{\alpha_n} \cup U_\alpha.
\]
So $\mathcal{C}_0 = \{U_{\alpha_1}, \dots, U_{\alpha_n}, U_\alpha\} \subset \mathcal{C}$ is a finite subcover of $[a,b]$, hence $[a,b]$ is compact.
\end{proof}

as this theorem, we can easily prove

\textbf{theorem:} if $F$ is closed set, then $F$ is compact set in compact space $(X,\mathcal{T})$

\begin{proposition}
A closed subset $F$ of a compact space $(X, \mathcal{T})$ is compact.
\end{proposition}

\begin{proof}
Let $C = \{U_\alpha \mid \alpha \in \Lambda, U_\alpha \in \mathcal{T}\}$ be an arbitrary open cover of $F$. Since $F \subset \bigcup_{\alpha \in \Lambda} U_\alpha$,
\[
C_X = C \cup \{F^c\} = \{U_\alpha, F^c \mid U_\alpha \in \mathcal{T}, \alpha \in \Lambda\}
\]
is an open cover of $X$ in $(X, \mathcal{T})$. Since $X$ is compact, there exists a finite subcover $C_X' \, (\subset C_X)$ of $X$, of the form
\[
\{U_{\alpha_i} \in C_X \mid i \in \{1, \dots, k\}\} \quad \text{or} \quad \{U_{\alpha_i} \in C_X \mid i \in \{1, \dots, k\}\} \cup \{F^c\}.
\]
Thus
\[
X \subset \bigcup_{i=1}^k U_{\alpha_i} \quad \text{or} \quad X \subset \left(\bigcup_{i=1}^k U_{\alpha_i}\right) \cup F^c.
\]
Since $F \cap F^c = \emptyset$, in either case $C_X'$ (excluding $F^c$ if present) is a finite subcover of $F$. Hence $F$ is compact in $(X, \mathcal{T})$.
\end{proof}

\textbf{limit point compact} is when topology $(X,\mathcal{T})$ has limit point on $X$

and this follows by compact space is limit point compact space. 












\end{document}
